\documentclass[twoside,english]{article}
\setlength{\oddsidemargin}{0.25 in}
\setlength{\evensidemargin}{-0.25 in}
\setlength{\topmargin}{-0.6 in}
\setlength{\textwidth}{6.5 in}
\setlength{\textheight}{8.5 in}
\setlength{\headsep}{0.75 in}
\setlength{\parindent}{0 in}
\setlength{\parskip}{0.03 in}

\usepackage{amsmath,amsfonts,graphicx,algorithm,caption}
\usepackage[noend]{algpseudocode}
\usepackage{hyperref}
\usepackage[numbers]{natbib}
\usepackage[document]{ragged2e}
\usepackage{mathtools}
\usepackage{xcolor}
\usepackage{booktabs}
\usepackage{enumitem}
\usepackage{multicol}
\usepackage{titlesec}
\titleformat*{\section}{\large\bfseries}
\titlespacing*{\section}{0pt}{6pt}{3pt}
\setlist{nosep,leftmargin=1.2em}
\captionsetup{font=footnotesize,skip=2pt}

\newcounter{lecnum}
\renewcommand{\thepage}{\arabic{page}}
\renewcommand{\thesection}{\arabic{section}}
\renewcommand{\theequation}{\arabic{equation}}
\renewcommand{\thefigure}{\arabic{figure}}
\renewcommand{\thetable}{\arabic{table}}

\newcommand{\lecture}[4]{
   \pagestyle{myheadings}
   \thispagestyle{plain}
   \newpage
   \setcounter{page}{1}
   \noindent
   \begin{center}
   \framebox{
      \vbox{\vspace{2mm}
    \hbox to 6.28in { {\bf IE 643: Deep Learning - Theory and Practice
		\hfill July-Dec 2026} }
       \vspace{2mm}
       \hbox to 6.28in { {\Large \hfill Project Report: Novelty Assessment \hfill} }
       \vspace{2mm}
       \hbox to 6.28in { {\large \hfill #1 \hfill} }
       \vspace{2mm}
       \hbox to 6.28in { {\it Team Name: #2 \hfill Team Members: #3} }
      \vspace{2mm}}
   }
   \end{center}
   \vspace*{1mm}
}
\renewcommand{\cite}[1]{[#1]}

\begin{document}
\lecture{Automated Respiratory Sound Classification via Audio Spectrogram Transformers}{doomed\_coders}{Soham Mhatre (26M2572), Parivat (26M2575)}

\textbf{Abstract.} We classify ICBHI 2017 respiratory cycles (normal, crackle, wheeze, both) with an Audio Spectrogram Transformer (AST) and ask whether generated spectrograms of the rare classes really help. Changing only the generated data, synthetic cycles add up to two points of ICBHI Score. Our best method, a GAN pretrained with the external HF\_Lung data, reaches an ensemble Score of \textbf{62.10}, up from 61.13 in our previous report.

\section{Introduction}
ICBHI 2017 is small and imbalanced, so we enlarge its rare classes with generated spectrograms. After the intensive assessment we first checked whether this helps at all: we trained the \emph{same} AST with and without WGAN-GP samples, keeping the data split, padding and training recipe fixed, and report sensitivity and specificity next to the Score. We then tested four ideas of our own, one at a time. All our GANs are BigGAN-style \cite{9} (hinge loss, spectral norm, class-conditional batch norm), not WGAN-GP, and real background noise from quiet frames of training cycles is mixed into every fake \cite{5}.
\begin{itemize}
\item \textbf{(1) Azimuthal integration} \cite{2}: the 2D Fourier power spectrum averaged over rings, used to find the GAN's spectral errors and as a training loss.
\item \textbf{(2) Up-sampling}: nearest-neighbour vs anti-aliased \cite{6} up-convolution, and our \emph{dual-branch generator}: a global CNN draws the overall spectrum and breathing envelope, a local CNN the fine detail, and a learned gate places the detail ($x=E+M\odot D$, cf.\ \cite{7}).
\item \textbf{(3) External data}: HF\_Lung\_V1 \cite{1} cycles (rhonchi-only ``wheeze'' cycles removed) train the generator together with ICBHI under a source label.
\item \textbf{(4) Diffusion}: a U-Net diffusion model (EDM \cite{3}) with classifier-free guidance \cite{4}; the scale $w$ sets how strongly a fake shows its class ($w=0$: not at all).
\end{itemize}

\section{Setup}
\textbf{Metric and classifier.} ICBHI Score $=(\mathrm{SE}+\mathrm{SP})/2$ on the official test split (SE, SP: recall of abnormal and normal cycles); 3 seeds per arm (5 where stated), reported as mean $\pm$ sd and as the ensemble of the seeds. The classifier is the same everywhere: an AudioSet-pretrained AST, fully fine-tuned, with a device-adversarial head, 8 s windows, 30 epochs, last epoch kept. Only the 3{,}000 generated training cycles (1{,}000 per rare class) change.

\textbf{Data.} ICBHI 2017: 920 recordings, 126 patients, 4 stethoscopes; 4{,}131 training cycles (2057 / 1210 / 501 / 363) and 2{,}756 test cycles. HF\_Lung\_V1 \cite{1}: 9{,}765 15-s recordings, 279 patients, 4 kHz, 27{,}168 labelled cycles, used to train generators only (plus one control).

\section{Results}

\textbf{Generated data helps once it is isolated.} On the same five seeds, GAN cycles raise the mean Score from 58.04 to 59.75 and diffusion cycles to 60.04 ($p=0.05$); equal-size ensembles rise from 60.00 to 61.73 and 62.00. With shuffled labels the gain disappears, so the class content of the fakes is what helps. WGAN-GP samples barely move the Score (+0.3).

\begin{table}[!ht]
\centering\footnotesize
\resizebox{\textwidth}{!}{\begin{tabular}{llcc}
\toprule
Comparison & Arm (only the generated data changes) & Mean $\pm$ sd (SE / SP) & Ensemble (SE / SP) \\
\midrule
WGAN-GP & AST, no generated data (previous report) & 59.39 $\pm$ 0.92 (45.0 / 73.8) & 61.13 (44.4 / 77.8) \\
WGAN-GP & + WGAN-GP samples & 59.71 $\pm$ 0.67 (46.4 / 73.1) & -- \\
\midrule
Ours & No generated data, 5 seeds & 58.04 $\pm$ 1.24 (45.7 / 70.4) & 60.00 (46.7 / 73.3) \\
Ours & + BigGAN-style GAN + real noise, 5 seeds & 59.75 $\pm$ 1.30 (45.3 / 74.2) & 61.73 (45.9 / 77.6) \\
Ours & + diffusion with HF ($w=4$), 5 seeds & \textbf{60.04 $\pm$ 1.00} (45.7 / 74.4) & \textbf{62.00} (45.8 / 78.2) \\
Ours & control: same GAN fakes, labels shuffled & 58.16 $\pm$ 0.83 (44.2 / 72.1) & 60.60 (44.7 / 76.5) \\
\bottomrule
\end{tabular}}
\caption{Same AST, split, padding and recipe; only the generated data changes. Ensembles have equal size (5 models).}
\end{table}

\begin{table}[!ht]
\centering\footnotesize
\resizebox{\textwidth}{!}{\begin{tabular}{lllcc}
\toprule
Idea & Arm (3 seeds) & Generator data & Mean $\pm$ sd (SE / SP) & Ensemble (SE / SP) \\
\midrule
(3) & GAN pretrained on clean HF, then ICBHI & ICBHI + HF & 59.55 $\pm$ 1.28 (45.5 / 73.6) & \textbf{62.10} (47.1 / 77.1) \\
(3)+(4) & Diffusion, $w=4$ & ICBHI + HF & \textbf{60.49 $\pm$ 0.25} (45.2 / 75.8) & 61.99 (46.0 / 78.0) \\
(4) & Diffusion, $w=2$ & ICBHI & 60.28 $\pm$ 1.45 (45.9 / 74.7) & 61.84 (46.0 / 77.7) \\
(1)+(2) & Dual-branch GAN + spectral loss & ICBHI & 60.01 $\pm$ 1.79 (45.1 / 75.0) & 61.75 (44.8 / 78.7) \\
(2) & GAN, nearest up-sampling + real noise & ICBHI & \textbf{60.50 $\pm$ 0.70} (44.5 / 76.5) & 61.53 (44.8 / 78.3) \\
(1) & GAN + Durall spectral loss & ICBHI & 59.45 $\pm$ 1.51 (45.1 / 73.8) & 60.59 (44.6 / 76.6) \\
(3) & HF as \emph{real} classifier data (control) & -- & 56.79 $\pm$ 1.14 (46.3 / 67.3) & 57.45 (45.6 / 69.3) \\
\bottomrule
\end{tabular}}
\caption{Each idea on its own, sorted by ensemble Score.}
\end{table}

\textbf{Azimuthal integration and up-sampling.} The 1D spectrum showed what was wrong with our first generator: nearest-neighbour up-convolution lost fine high-frequency detail. Noise injection and diffusion both shrank this gap about five-fold. For the classifier, the most useful fix was mixing real stethoscope noise into every fake: the GAN trained on ICBHI alone then became our best single model, 60.50 $\pm$ 0.70. The azimuthal spectrum as a loss did not help (59.45). The dual-branch generator worked only once its gate was kept from collapsing, reaching an ensemble of 61.75 with the spectral loss.

\textbf{External data.} HF\_Lung helps only through the generator. Pretraining the GAN on clean HF cycles raises the ensemble from 61.53 (ICBHI only) to \textbf{62.10}, our best method, and in diffusion HF makes results much more stable (sd 0.25 vs 1.45). Given to the classifier directly, HF lowers the Score to 56.79: its patients, devices and labels differ from ICBHI.

\textbf{Diffusion.} On ICBHI alone diffusion reaches 60.28 (ensemble 61.84). Guidance matters: $w=0$ gives 58.58, $w=4$ 60.49, stronger guidance adds nothing. Most of the gain comes on the Littmann 3200 stethoscope, which has only 41 training cycles (46 $\rightarrow$ 52--54).

\section{Conclusion}
Generated spectrograms do help when everything else is fixed: about +1.7 points for the GAN and +2.0 for diffusion. Three things made the difference: mixing real stethoscope noise into the fakes, using HF\_Lung to train the generator rather than the classifier, and moderate classifier-free guidance. Our best method is the GAN pretrained with HF\_Lung (ensemble 62.10); our best single model is the improved-up-sampling GAN on ICBHI only (60.50 $\pm$ 0.70). Pooling all 20 diffusion models gives 62.54, the highest Score we reached, but as a large ensemble rather than a practical method. The guidance scale was compared on test data, so diffusion numbers are slightly optimistic. Variants that did not help, against our best GAN (60.50): anti-aliased up-sampling \cite{6} 58.14 ($-$2.36), flow matching \cite{8} 58.74 ($-$1.76), an HF event-detector teacher 59.35 ($-$1.15), and rule-based event insertion 56.42 ($-$4.08).

\vspace{-1mm}
\renewcommand{\refname}{\normalsize References}
{\scriptsize
\setlength{\columnsep}{10pt}\begin{multicols}{2}
\begin{thebibliography}{99}\setlength{\itemsep}{0pt}\setlength{\parskip}{0pt}
\bibitem{1} Fu-Shun Hsu, Shang-Ran Huang, Chien-Wen Huang, et al. Benchmarking of eight recurrent neural network variants for breath phase and adventitious sound detection on a self-developed open-access lung sound database -- HF\_Lung\_V1, PLoS ONE, 2021. https://arxiv.org/abs/2102.03049
\bibitem{2} Ricard Durall, Margret Keuper and Janis Keuper. Watch your Up-Convolution: CNN Based Generative Deep Neural Networks are Failing to Reproduce Spectral Distributions, CVPR, 2020. https://arxiv.org/abs/2003.01826
\bibitem{3} Tero Karras, Miika Aittala, Timo Aila and Samuli Laine. Elucidating the Design Space of Diffusion-Based Generative Models, NeurIPS, 2022. https://arxiv.org/abs/2206.00364
\bibitem{4} Jonathan Ho and Tim Salimans. Classifier-Free Diffusion Guidance, NeurIPS Workshop, 2021. https://arxiv.org/abs/2207.12598
\bibitem{5} Justin Salamon and Juan Pablo Bello. Deep Convolutional Neural Networks and Data Augmentation for Environmental Sound Classification, IEEE Signal Processing Letters, 2017. https://arxiv.org/abs/1608.04363
\bibitem{6} Richard Zhang. Making Convolutional Networks Shift-Invariant Again, ICML, 2019. https://arxiv.org/abs/1904.11486
\bibitem{7} Jianwei Yang, Anitha Kannan, Dhruv Batra and Devi Parikh. LR-GAN: Layered Recursive Generative Adversarial Networks for Image Generation, ICLR, 2017. https://arxiv.org/abs/1703.01560
\bibitem{8} Yaron Lipman, Ricky T. Q. Chen, Heli Ben-Hamu, Maximilian Nickel and Matt Le. Flow Matching for Generative Modeling, ICLR, 2023. https://arxiv.org/abs/2210.02747
\bibitem{9} Andrew Brock, Jeff Donahue and Karen Simonyan. Large Scale GAN Training for High Fidelity Natural Image Synthesis, ICLR, 2019. https://arxiv.org/abs/1809.11096
\end{thebibliography}
\end{multicols}}
\end{document}
